% =============================================================================
% Shared preamble for update_181126: the architecture deck.
% Visual vocabulary follows the two references in this folder:
%   image copy.png  -> 1.1 (function columns, colour-coded data lines)
%   image.png       -> sub-component slides (one stage panel: picture + caption)
% =============================================================================
\usetheme{default}
\usecolortheme{seagull}
\setbeamertemplate{navigation symbols}{}
\setbeamertemplate{footline}[frame number]
\setbeamerfont{frametitle}{size=\large,series=\bfseries}
\setbeamercolor{frametitle}{fg=black!85}
\usepackage[T1]{fontenc}
\usepackage[utf8]{inputenc}
\usepackage{booktabs}
\usepackage{array}
\usepackage{multirow}
\usepackage{siunitx}
\usepackage{graphicx}
\usepackage{xcolor}
\usepackage{amssymb}
\usepackage{adjustbox}
\usepackage[most]{tcolorbox}
\usepackage{tikz}
\usetikzlibrary{arrows.meta,positioning,fit,calc,backgrounds,
                decorations.pathreplacing,shapes.geometric,shadows}
\sisetup{detect-all,range-phrase=--,range-units=single,retain-explicit-plus}
\DeclareSIUnit{\mebi}{Mi}
\newcommand{\repo}[1]{\texttt{\detokenize{#1}}}
\newcommand{\src}[1]{\vfill{\tiny\color{black!50}#1}}

% ---------------------------------------------------------------- colours
% function columns (1.1), after image copy.png
\definecolor{sensorcol}{HTML}{2FB5E8}
\definecolor{perccol}{HTML}{26A69A}
\definecolor{plancol}{HTML}{F2735A}
\definecolor{ctrlcol}{HTML}{8FA3BA}
\definecolor{mapcol}{HTML}{34495E}
% one colour per system
\definecolor{orbc}{HTML}{D6336C}
\definecolor{vinsc}{HTML}{E67E22}
\definecolor{rtabc}{HTML}{27AE60}
\definecolor{navc}{HTML}{2E86DE}
\definecolor{sysc}{HTML}{5D6D7E}
% data-line colours: one per kind of data
\definecolor{dimg}{HTML}{E74C3C}
\definecolor{dimu}{HTML}{F39C12}
\definecolor{dwhl}{HTML}{16A085}
\definecolor{dpose}{HTML}{8E44AD}
\definecolor{dmap}{HTML}{2C3E50}
\definecolor{dcmd}{HTML}{2E86DE}
% marks
\definecolor{ourc}{HTML}{E0A100}
\definecolor{hlc}{HTML}{C0392B}
\definecolor{offc}{HTML}{AAB7B8}
\definecolor{good}{HTML}{1E8449}
\definecolor{bad}{HTML}{C0392B}

% ---------------------------------------------------------------- tikz styles
\tikzset{
  box/.style={draw=black!45, rounded corners=2pt, align=center, font=\scriptsize,
              inner sep=2pt, minimum height=8mm, text width=22mm, fill=white},
  arr/.style={-{Latex[length=1.8mm]}, thick, black!65},
  darr/.style={arr, dashed},
  rate/.style={font=\tiny, fill=white, inner sep=1pt, text=black!75, align=center},
  % "added by this project" and "switched off"
  ours/.style={draw=ourc, dashed, line width=1.2pt},
  off/.style={draw=offc, dashed, fill=offc!15, text=black!45},
  % the current component
  hl/.style={draw=hlc, line width=1.6pt, fill=yellow!35},
  % locator chips
  chip/.style={draw=black!35, rounded corners=1.5pt, align=center, font=\scriptsize,
               inner sep=1.5pt, minimum height=8mm, text width=17mm},
  grouplbl/.style={font=\scriptsize\bfseries, text=black!60},
  brace/.style={decorate, decoration={brace, amplitude=3pt, raise=2pt}, thick, black!40},
  % small drawing helpers for the stage panels
  pt/.style={circle, fill=#1, inner sep=0pt, minimum size=1.6mm},
  pt/.default=black,
  cam/.style={isosceles triangle, shape border rotate=#1, draw=black!70, fill=white,
              minimum width=4mm, minimum height=4mm, inner sep=0pt},
  cam/.default=0,
}
% Every component has a style of its own name, empty by default.
% \hl{ORB3} inside a frame turns that one component yellow-and-red, so the
% locator strip and the block diagram both show "you are here".
\tikzset{ORB1/.style={},ORB2/.style={},ORB3/.style={},ORB4/.style={},
         ORB5/.style={},ORB6/.style={},ORB7/.style={},ORB8/.style={},
         VINS1/.style={},VINS2/.style={},VINS3/.style={},VINS4/.style={},
         VINS5/.style={},VINS6/.style={},VINS7/.style={},
         RTAB1/.style={},RTAB2/.style={},RTAB3/.style={},RTAB4/.style={},
         RTAB5/.style={},RTAB6/.style={},RTAB7/.style={},RTAB8/.style={},RTAB9/.style={},
         NAV1/.style={},NAV2/.style={},NAV3/.style={},NAV4/.style={},
         NAV5/.style={},NAV6/.style={},NAV7/.style={},NAV8/.style={}}
\newcommand{\hl}[1]{\tikzset{#1/.style={hl}}}

% ---------------------------------------------------------------- panels
% The stage panel: a coloured header, one picture, one caption.  The picture is
% scaled to fit, never to overflow.
\newtcolorbox{stagepanel}[2][]{enhanced, colback=white, colframe=#2,
  boxrule=0.8pt, arc=2pt, left=3pt, right=3pt, top=2pt, bottom=2pt,
  fonttitle=\scriptsize\bfseries, coltitle=white, title={#1}}
\newcommand{\stage}[4]{% #1 colour, #2 header, #3 tikz body, #4 caption
  \begin{stagepanel}[#2]{#1}
    \centering
    \adjustbox{max width=\linewidth, max height=3.3cm}{#3}\par\vspace{2pt}
    {\scriptsize #4}
  \end{stagepanel}}
% The right-hand column of a component slide.
\newtcolorbox{linkpanel}{enhanced, colback=mapcol!5, colframe=mapcol,
  boxrule=0.8pt, arc=2pt, left=3pt, right=3pt, top=2pt, bottom=2pt,
  fonttitle=\scriptsize\bfseries, coltitle=white, title={LINK TO RESULTS}}
\newcommand{\side}[5]{% #1 why, #2 in, #3 out, #4 can go wrong, #5 signal
  {\scriptsize\textbf{\color{mapcol}WHY IT EXISTS}\par\footnotesize #1\par}
  \vspace{0.3em}
  {\scriptsize\raggedright\textbf{\color{mapcol}IN}\quad #2\par
   \textbf{\color{mapcol}OUT}\enspace #3\par}
  \vspace{0.3em}
  \begin{linkpanel}\scriptsize\raggedright
    \textbf{Can go wrong:} #4\par\smallskip
    \textbf{Signal:} #5\par\smallskip
    \textbf{Results:} \resref
  \end{linkpanel}}
% A whole component slide body: strip, then panel (left) + side (right).
\newcommand{\compbody}[9]{% #1 strip macro, #2 colour, #3 tikz, #4 caption,
                          % #5 why, #6 in, #7 out, #8 wrong, #9 signal
  #1\par\vspace{0.25em}
  \begin{columns}[T]
    \column{0.55\textwidth}\stage{#2}{HOW IT WORKS}{#3}{#4}
    \column{0.43\textwidth}\side{#5}{#6}{#7}{#8}{#9}
  \end{columns}}

% One lane of the "three clocks" timeline (1.1): y, label, rate in Hz, colour.
\newcommand{\lane}[4]{%
  \node[anchor=east, align=right] at (-0.02,#1) {#2};
  \draw[black!15] (0,#1) -- (1,#1);
  \foreach \k in {1,...,#3} \draw[#4, line width=0.7pt]
    ({(\k-0.5)/#3},#1-0.18) -- ({(\k-0.5)/#3},#1+0.18);}

% ---------------------------------------------------------------- results
% A component slide's "Results" line.  \res{...} inside a frame sets it for that
% frame only; \R{n}{text} links to the finding slide labelled Rn.
\newcommand{\resref}{{\color{black!40}\emph{no finding yet}}}
\newcommand{\res}[1]{\renewcommand{\resref}{#1}}
\newcommand{\R}[2]{\hyperlink{R#1}{\textbf{\color{hlc}R#1}}~#2}
\graphicspath{{../../final_report/figures/}{../experiment_4/fig/}}
% A finding slide body, under its strip(s): picture left, two panels right.
\newcommand{\finding}[5]{% #1 colour, #2 picture, #3 what we saw, #4 why, #5 caveat
  \begin{columns}[T]
    \column{0.57\textwidth}\centering
      \adjustbox{max width=\linewidth, max height=0.50\textheight}{#2}
    \column{0.42\textwidth}
      \begin{stagepanel}[WHAT WE SAW]{mapcol}\scriptsize\raggedright #3\end{stagepanel}
      \vspace{-0.4em}
      \begin{stagepanel}[WHY: INSIDE THE BLOCK]{#1}\scriptsize\raggedright #4\end{stagepanel}
  \end{columns}
  \vspace{0.2em}
  {\scriptsize\color{black!70}\textbf{\color{bad}Careful:} #5\par}}
